% GENERATED FILE. Edit canonical lessons, then run npm run generate:books.
% canonical-source-hash: fnv1a64:a8225ddc84468830
% canonical-lessons: HI-C512-kapas, HI-C512-carkha, HI-C512-kaatna, HI-C512-kargha, HI-C512-gath

\chapter{Thread and Loom}
\label{ch:hi-thread-and-loom}
\hlchaptermodality{512}

\begin{chapteropening}
\textbf{By the end of this chapter:} \emph{I can talk about cotton, spinning and weaving.}

Complete the last lesson of chapter 512: I can talk about cotton, spinning and weaving.
\end{chapteropening}

\section[\texorpdfstring{\hi{कपास} (kapās)}{कपास (kapās)}]{\hi{कपास} (kapās) — cotton (the plant, the raw fibre)}

\label{lesson:HI-C512-kapas}



\begin{quote}
Before the new one: say the Hindi for a mould, then the Hindi for to pour into a mould.
\end{quote}

\subsection*{You'll want to know: \hi{कपास}}
\textbf{\hi{कपास}} — \emph{kapās} — \textquotedblleft{}cotton (the plant, the raw fibre)\textquotedblright{}.

\begin{quote}
\hi{किसान} \hi{ने} \hi{खेत} \hi{में} \hi{कपास} \hi{बोई।}
\end{quote}

\begin{quote}
\emph{Kisān ne khet meṁ kapās boī.}
\end{quote}

\begin{quote}
The farmer sowed cotton in the field.
\end{quote}

Feminine.

\begin{cousinweb}
\textbf{\hi{कपास}} (\emph{kapās}) is a tadbhava, inherited from Sanskrit \emph{karpāsa}, \textquotedblleft{}the cotton plant\textquotedblright{} (CDIAL 2877).
\end{cousinweb}

\begin{grammarlens}[title={What you've built}]
One more word for describing what you lived and felt.
\end{grammarlens}

\subsection*{Your turn}
\begin{itemize}
\raggedright
  \item \emph{Say it:} \emph{kapās}
  \item \emph{Say it:} \emph{kapās}, once more
  \item \emph{Say it:} \emph{Kisān ne khet meṁ kapās boī.} — The farmer sowed cotton in the field.
  \item \emph{Recall:} say the Hindi for a mould, then the Hindi for to pour into a mould, then say \emph{kapās} again
\end{itemize}

\begin{tcolorbox}[breakable,colback=teal!4,colframe=teal!35!black,title={Before you move on}]
What does \hi{कपास} mean? (\textquotedblleft{}Cotton (the plant, the raw fibre)\textquotedblright{}.) Say it once more.
\end{tcolorbox}
\section[\texorpdfstring{\hi{चरखा} (carkhā)}{चरखा (carkhā)}]{\hi{चरखा} (carkhā) — a spinning wheel}

\label{lesson:HI-C512-carkha}



\begin{quote}
Before the new one: say the Hindi for to pour into a mould, then the Hindi for cotton (the plant, the raw fibre).
\end{quote}

\subsection*{You'll want to know: \hi{चरखा}}
\textbf{\hi{चरखा}} — \emph{carkhā} — \textquotedblleft{}a spinning wheel\textquotedblright{}.

\begin{quote}
\hi{मेरी} \hi{नानी} \hi{हर} \hi{शाम} \hi{चरखा} \hi{चलाती} \hi{थी।}
\end{quote}

\begin{quote}
\emph{Merī nānī har śām carkhā calātī thī.}
\end{quote}

\begin{quote}
My grandmother used to work the spinning wheel every evening.
\end{quote}

Masculine.

\begin{cousinweb}
\textbf{\hi{चरखा}} (\emph{carkhā}) is usually traced to Persian \emph{charkha}, \textquotedblleft{}a wheel\textquotedblright{}.
\end{cousinweb}

\begin{grammarlens}[title={What you've built}]
One more word for describing what you lived and felt.
\end{grammarlens}

\subsection*{Your turn}
\begin{itemize}
\raggedright
  \item \emph{Say it:} \emph{carkhā}
  \item \emph{Say it:} \emph{carkhā}, once more
  \item \emph{Say it:} \emph{Merī nānī har śām carkhā calātī thī.} — My grandmother used to work the spinning wheel every evening.
  \item \emph{Recall:} say the Hindi for to pour into a mould, then the Hindi for cotton (the plant, the raw fibre), then say \emph{carkhā} again
\end{itemize}

\begin{tcolorbox}[breakable,colback=teal!4,colframe=teal!35!black,title={Before you move on}]
What does \hi{चरखा} mean? (\textquotedblleft{}A spinning wheel\textquotedblright{}.) Say it once more.
\end{tcolorbox}
\section[\texorpdfstring{\hi{कातना} (kātnā)}{कातना (kātnā)}]{\hi{कातना} (kātnā) — to spin (thread)}

\label{lesson:HI-C512-kaatna}



\begin{quote}
Before the new one: say the Hindi for cotton (the plant, the raw fibre), then the Hindi for a spinning wheel.
\end{quote}

\subsection*{You'll want to know: \hi{कातना}}
\textbf{\hi{कातना}} — \emph{kātnā} — \textquotedblleft{}to spin (thread)\textquotedblright{}.

\begin{quote}
\hi{वह} \hi{चरखे} \hi{पर} \hi{कपास} \hi{से} \hi{धागा} \hi{कात} \hi{रही} \hi{है।}
\end{quote}

\begin{quote}
\emph{Vah carkhe par kapās se dhāgā kāt rahī hai.}
\end{quote}

\begin{quote}
She is spinning thread from cotton on the wheel.
\end{quote}

Past, with \hi{ने} (he or she …): \hi{उसने} \hi{धागा} \hi{काता}, \emph{usne dhāgā kātā}.

\begin{cousinweb}
\textbf{\hi{कातना}} (\emph{kātnā}) is a tadbhava, inherited from \emph{*kartati}, \textquotedblleft{}spins\textquotedblright{}, a form CDIAL reconstructs from its descendants (CDIAL 2855).
\end{cousinweb}

\begin{grammarlens}[title={What you've built}]
One more word for describing what you lived and felt.
\end{grammarlens}

\subsection*{Your turn}
\begin{itemize}
\raggedright
  \item \emph{Say it:} \emph{kātnā}
  \item \emph{Say it:} \emph{kātnā}, once more
  \item \emph{Say it:} \emph{Vah carkhe par kapās se dhāgā kāt rahī hai.} — She is spinning thread from cotton on the wheel.
  \item \emph{Recall:} say the Hindi for cotton (the plant, the raw fibre), then the Hindi for a spinning wheel, then say \emph{kātnā} again
\end{itemize}

\begin{tcolorbox}[breakable,colback=teal!4,colframe=teal!35!black,title={Before you move on}]
What does \hi{कातना} mean? (\textquotedblleft{}To spin (thread)\textquotedblright{}.) Say it once more.
\end{tcolorbox}
\section[\texorpdfstring{\hi{करघा} (karghā)}{करघा (karghā)}]{\hi{करघा} (karghā) — a loom}

\label{lesson:HI-C512-kargha}



\begin{quote}
Before the new one: say the Hindi for a spinning wheel, then the Hindi for to spin (thread).
\end{quote}

\subsection*{You'll want to know: \hi{करघा}}
\textbf{\hi{करघा}} — \emph{karghā} — \textquotedblleft{}a loom\textquotedblright{}.

\begin{quote}
\hi{करघे} \hi{पर} \hi{वह} \hi{एक} \hi{दरी} \hi{बुन} \hi{रहा} \hi{है।}
\end{quote}

\begin{quote}
\emph{Karghe par vah ek darī bun rahā hai.}
\end{quote}

\begin{quote}
He is weaving a rug on the loom.
\end{quote}

Masculine.

\begin{cousinweb}
\textbf{\hi{करघा}} (\emph{karghā}) is usually traced to Persian \emph{kārgāh}, \textquotedblleft{}a workplace\textquotedblright{}. Hindi narrowed it to the weaver's loom.
\end{cousinweb}

\begin{grammarlens}[title={What you've built}]
One more word for describing what you lived and felt.
\end{grammarlens}

\subsection*{Your turn}
\begin{itemize}
\raggedright
  \item \emph{Say it:} \emph{karghā}
  \item \emph{Say it:} \emph{karghā}, once more
  \item \emph{Say it:} \emph{Karghe par vah ek darī bun rahā hai.} — He is weaving a rug on the loom.
  \item \emph{Recall:} say the Hindi for a spinning wheel, then the Hindi for to spin (thread), then say \emph{karghā} again
\end{itemize}

\begin{tcolorbox}[breakable,colback=teal!4,colframe=teal!35!black,title={Before you move on}]
What does \hi{करघा} mean? (\textquotedblleft{}A loom\textquotedblright{}.) Say it once more.
\end{tcolorbox}
\section[\texorpdfstring{\hi{गाँठ} (gā̃ṭh)}{गाँठ (gā̃ṭh)}]{\hi{गाँठ} (gā̃ṭh) — a knot}

\label{lesson:HI-C512-gath}



\begin{quote}
Before the new one: say the Hindi for to spin (thread), then the Hindi for a loom.
\end{quote}

\subsection*{You'll want to know: \hi{गाँठ}}
\textbf{\hi{गाँठ}} — \emph{gā̃ṭh} — \textquotedblleft{}a knot\textquotedblright{}.

\begin{quote}
\hi{धागे} \hi{में} \hi{एक} \hi{गाँठ} \hi{पड़} \hi{गई।}
\end{quote}

\begin{quote}
\emph{Dhāge meṁ ek gā̃ṭh paṛ gaī.}
\end{quote}

\begin{quote}
A knot has formed in the thread.
\end{quote}

Feminine.

\begin{cousinweb}
\textbf{\hi{गाँठ}} (\emph{gā̃ṭh}) is a tadbhava, inherited from Sanskrit \emph{granthi}, \textquotedblleft{}a knot\textquotedblright{} (CDIAL 4354).
\end{cousinweb}

\begin{grammarlens}[title={What you've built}]
One more word for describing what you lived and felt.
\end{grammarlens}

\subsection*{Your turn}
\begin{itemize}
\raggedright
  \item \emph{Say it:} \emph{gā̃ṭh}
  \item \emph{Say it:} \emph{gā̃ṭh}, once more
  \item \emph{Say it:} \emph{Dhāge meṁ ek gā̃ṭh paṛ gaī.} — A knot has formed in the thread.
  \item \emph{Recall:} say the Hindi for to spin (thread), then the Hindi for a loom, then say \emph{gā̃ṭh} again
\end{itemize}

\begin{tcolorbox}[breakable,colback=teal!4,colframe=teal!35!black,title={Before you move on}]
What does \hi{गाँठ} mean? (\textquotedblleft{}A knot\textquotedblright{}.) Say it once more.
\end{tcolorbox}
